\verseline{}{与此同时，全能的奥林匹斯宫门洞开，}
\verseline{}{众神之父、人类之王召集神祇}
\verseline{}{来到星光灿烂的殿堂；他从高处俯临，}
\verseline{}{望遍大地、达耳达尼亚人的营垒与拉丁各族。\footnote{奥林匹斯的全能者之宫洞开，朱庇特以“众神之父、人类之王”的称号召集会议。他从天上俯瞰大地、达耳达尼亚人的营垒与拉丁各族；开篇由神廷转到凡间战场。}}
\verseline{5}{众神在双门敞开的宫中坐定，他先开口：}
\verseline{}{“天上的大神们，为何你们的意志}
\verseline{}{竟全然倒转，怀着偏私的心彼此争执？}
\verseline{}{我曾禁止意大利人与特乌克罗斯人交战。}
\verseline{}{是谁违背禁令挑起纷争？何种恐惧怂恿}
\verseline{10}{这一方或那一方追逐兵戈、以剑挑衅？}
\verseline{}{正当的交战之日自会来到，不要催它降临；}
\verseline{}{终有一日，凶悍的迦太基将越过}
\verseline{}{敞开的阿尔卑斯山，把浩劫送向罗马城堡：}
\verseline{}{到那时，才可用仇恨相争，才可劫掠财物。\footnote{朱庇特所说的未来战争预指汉尼拔越过阿尔卑斯山、威胁罗马的第二次布匿战争。神明禁止双方此刻把仇恨倾泻到底，却说罗马人与迦太基人的大战应留到后世；眼前的意大利战争因此同未来的布匿战争并提。}}
\verseline{15}{如今住手吧，欣然订立双方认可的盟约。”}
\verseline{}{朱庇特只说了这些；金色的维纳斯却以}
\verseline{}{更长的话回答：}
\verseline{}{“父亲啊，人类与万物的永恒主宰！}
\verseline{}{事已至此，我们还能向谁求告？}
\verseline{20}{你可看见鲁图利人怎样凌辱我们，图尔努斯怎样}
\verseline{}{骑着名马，骄然穿过众军，乘战势}
\verseline{}{冲杀吗？紧闭的城墙已护不住特乌克罗斯人：}
\verseline{}{敌人竟在门内、就在城垣土垒上}
\verseline{}{与他们交战，让壕沟涨满鲜血。}
\verseline{25}{埃涅阿斯毫不知情，远在他方。你永远不肯}
\verseline{}{解除围困吗？敌人又一次逼近}
\verseline{}{新生特洛伊的城墙，另一支军队也已到来；}
\verseline{}{提丢斯之子又从埃托利亚的阿耳皮起兵，}
\verseline{}{再度攻向特乌克罗斯人。我相信，留给我的仍是旧伤，}
\verseline{30}{我虽是你的女儿，还须等待凡人的刀枪。\footnote{“提丢斯之子”是狄俄墨得斯。他在《伊利亚特》第五卷中曾刺伤维纳斯；阿耳皮是传说中他流落意大利后建立的城市。维纳斯借旧伤预想自己将再次面对凡人的兵器。“又起兵”是她辩词中的威胁说法；第十一卷使者回报，狄俄墨得斯拒绝参战。}}
\verseline{}{若特洛伊人未经你准许、违逆你的神意}
\verseline{}{来到意大利，就让他们偿还罪过，不要}
\verseline{}{用援助扶持他们；但若他们遵从天神与}
\verseline{}{亡魂赐下的众多神谕，为何如今竟有人}
\verseline{35}{能够推翻你的命令，另立新的命运？}
\verseline{}{何必再提厄律克斯海岸被焚的舰队，}
\verseline{}{再提风暴之王和从埃俄利亚唤起的}
\verseline{}{狂风，或受云驱使的伊里斯？}
\verseline{}{如今她甚至惊动亡魂——这是世事中}
\verseline{40}{唯一尚未尝试的一着——突然把阿勒克托}
\verseline{}{放上人间，使她狂舞着穿过意大利诸城。\footnote{维纳斯依次追忆此前各卷中朱诺一方的干预：厄律克斯海岸的焚船、埃俄卢斯放出的狂风、伊里斯传令，以及第七卷中阿勒克托被遣往意大利诸城。她是在神界争辩中陈述己方所受侵害，所列责任集中于朱诺一方。}}
\verseline{}{我已不再为统治权忧心；时势尚在时，}
\verseline{}{我们曾有过这份希望。让你愿其胜者取胜吧。}
\verseline{}{倘若你那严酷的妻子不肯给特乌克罗斯人}
\verseline{45}{任何一方土地，父亲，我指着覆亡的特洛伊}
\verseline{}{仍在冒烟的废墟恳求：请准许}
\verseline{}{阿斯卡尼乌斯安然离开战阵，让我的孙儿活下去。}
\verseline{}{至于埃涅阿斯，尽可被抛荡在陌生的海上，}
\verseline{}{沿命运为他给出的任何道路而去：}
\verseline{50}{只让我能够庇护这孩子，带他脱离惨烈的战争。}
\verseline{}{我有阿玛图斯，有高耸的帕福斯和库忒拉，}
\verseline{}{也有伊达利乌姆的宫宇；让他放下兵器、无名无誉，}
\verseline{}{在那里度过一生。你尽可下令强大的}
\verseline{}{迦太基统治奥索尼亚；推罗人的城市不会再受阻挠。\footnote{阿玛图斯、帕福斯、库忒拉和伊达利乌姆均与维纳斯崇拜相关。inglorius 修饰想象中放下武器的阿斯卡尼乌斯，意为不再追求或获得战争荣名。维纳斯以放弃罗马未来乃至容许迦太基统治意大利，逼显恳求的极端性。}}
\verseline{55}{既然特乌克罗斯人寻找拉提乌姆与复生的帕加马，}
\verseline{}{逃过战争的灾祸，穿过阿尔戈斯的烈火，}
\verseline{}{又历尽海洋和辽阔陆地的重重险难，}
\verseline{}{这一切究竟有何益处？}
\verseline{}{留坐在祖国最后的灰烬上、}
\verseline{60}{留在特洛伊曾经所在的故土岂不更好？请把}
\verseline{}{克珊托斯河与西摩伊斯河还给这些苦难者吧，父亲，}
\verseline{}{让特乌克罗斯人重新经历特洛伊的灾厄。”}
\verseline{}{王后朱诺被沉重的怒火驱动：“为何逼我}
\verseline{}{打破深深的沉默，用言语揭开封闭的痛楚？}
\verseline{65}{凡人或神明中，有谁强迫埃涅阿斯}
\verseline{}{追逐战争，或以敌人的身份向拉提努斯王发动攻击？}
\verseline{}{就算他受命运指引来到意大利，}
\verseline{}{又受卡珊德拉的狂言催逼：难道是我们劝他离开营垒，}
\verseline{}{或把生命交给风浪？}
\verseline{70}{难道是我们劝他把战争大权和营墙托给一个孩子，}
\verseline{}{又去寻求伊特鲁里亚人的盟援，或惊扰安居的民族？}
\verseline{}{哪一位神、我们的哪一种残酷力量把他}
\verseline{}{诱入这场灾祸？朱诺在哪里，从云中降下的伊里斯又在哪里？}
\verseline{}{意大利人用火包围新生的特洛伊，便算不义，}
\verseline{75}{而图尔努斯立足自己的国土也算不义——}
\verseline{}{他的祖父是皮卢姆努斯，母亲是女神维尼利亚；}
\verseline{}{那么特洛伊人用黑色火炬向拉丁人施暴，}
\verseline{}{以轭压服异乡田野，掠走战利品，又算什么？}
\verseline{}{他们自选岳父，从怀中夺走许给别人的新娘，}
\verseline{80}{手持橄榄枝求和，却在船尾挂出兵器，又算什么？}
\verseline{}{你能从希腊人手中救出埃涅阿斯，}
\verseline{}{以云雾和虚幻的风替换这个男子，}
\verseline{}{也能把整支舰队变成同等数目的宁芙；}
\verseline{}{而我们稍稍援助鲁图利人，竟成了神律所禁？}
\verseline{85}{埃涅阿斯毫不知情地远去——就让他毫不知情地远去。}
\verseline{}{你有帕福斯与伊达利乌姆，也有高耸的库忒拉；}
\verseline{}{为何还试探这座孕育战争的城市与这些强悍的心？}
\verseline{}{试图从根基推翻佛里吉亚残破国运的，}
\verseline{}{究竟是我们，还是把苦难的特洛伊人交给阿凯亚人的那一方？}
\verseline{90}{是什么原因使欧罗巴、亚细亚一齐武装，}
\verseline{}{又是谁以窃夺毁坏盟约？}
\verseline{}{达耳达尼亚的奸夫难道在我引领下攻破斯巴达？}
\verseline{}{是我给他兵器，用欲望助长战争吗？}
\verseline{}{那时你才应为你的人恐惧；如今你来得太晚，}
\verseline{95}{只抛出不实的怨诉和无效的争辩。”\footnote{朱诺的答辩同样带有辩论立场。她把埃涅阿斯的离营、伊特鲁里亚联盟和特洛伊人争婚写成主动侵扰，又把特洛伊战争的责任推回帕里斯掳走海伦； “达耳达尼亚的奸夫”即帕里斯。她与维纳斯各自选取有利于己方的事实相互辩驳。}}
\verseline{}{朱诺如此陈说，天上的神祇全都}
\verseline{}{以各异的意见喧响，仿佛最初的风息}
\verseline{}{被困在林间，翻卷幽暗的低鸣，}
\verseline{}{向水手预告将来的风暴。}
\verseline{100}{于是掌有万物最高权力的全能之父}
\verseline{}{开口；当他说话，众神高大的宫宇沉寂，}
\verseline{}{大地从根基震颤，高远的天空无声，}
\verseline{}{西风停歇，海也以平静压住水面：\footnote{马克罗比乌斯（6.2.26）把这一整幅宇宙静默与恩尼乌斯《西庇阿》残篇并列：在那里天空静止，海神停息，太阳勒住马车，河流与树木也一同止息。维吉尔把这一静默置于朱庇特裁决诸神争执之前。}}
\verseline{}{“那么用心听取，把我的这些话钉牢在心中。}
\verseline{105}{既然不容奥索尼亚人与特乌克罗斯人缔盟，}
\verseline{}{你们的争执也没有终结，}
\verseline{}{今日每个人有怎样的时运、追随怎样的希望，}
\verseline{}{无论是特洛伊人还是鲁图利人，我都一视同仁——}
\verseline{}{不论特洛伊人的营垒是因意大利人的命运而遭到围困，}
\verseline{110}{还是特洛伊人误入歧途、听信不祥神示，皆然。}
\verseline{}{我也不为鲁图利人开脱：各自开端的事业}
\verseline{}{会给各自带来辛劳与命运；朱庇特为王，对众生相同。}
\verseline{}{命运自会找到道路。”他随即以冥府兄弟的斯提克斯河、\footnote{朱庇特宣布对特洛伊人与鲁图利人“不作区分”，限定于此刻不依阵营分配胜负；全诗既定的命运仍将自行实现。Fata viam invenient 指命运自会找到实现自身的道路。冥河誓言是神明最重的誓言；朱庇特点头使奥林匹斯震动。第一卷他向维纳斯确认埃涅阿斯后裔的未来时，也以点头作保证。}}
\verseline{}{以翻涌沥青、黑涡旋转的河岸为誓，}
\verseline{115}{点头应允；这一点头使整个奥林匹斯震颤。}
\verseline{}{言语至此终止。朱庇特从金色宝座上}
\verseline{}{起身，天神簇拥着他走向宫门。}
\verseline{}{与此同时，鲁图利人从四面紧逼营门，}
\verseline{}{要以杀戮扫倒守军，以火围住城墙。}
\verseline{120}{埃涅阿斯的军团却被围困在壁垒内，}
\verseline{}{毫无逃脱的希望。苦难的人徒然立在高塔，}
\verseline{}{以稀疏的人环守城垣：}
\verseline{}{伊姆布拉索斯之子阿西乌斯、希刻塔翁之子提摩特斯，}
\verseline{}{两个阿萨剌科斯和年长的提姆布里斯与卡斯托耳，}
\verseline{125}{组成前列；萨尔佩冬的两个兄弟}
\verseline{}{克拉鲁斯与泰蒙，也从高远的吕喀亚随他们守卫。}
\verseline{}{利耳涅索斯的阿克蒙用全身力量托举巨岩，}
\verseline{}{那是山体不小的一角；他的力量不逊于父亲}
\verseline{}{克吕提乌斯，也不逊于兄弟门涅斯透斯。}
\verseline{130}{有的用投枪，有的用石块竭力守御，}
\verseline{}{有的运来火焰，有的把箭搭上弓弦。}
\verseline{}{看哪，在众人之间，维纳斯最有理由牵挂的}
\verseline{}{达耳达尼亚少年露出端正的头颅，没有戴盔，}
\verseline{}{犹如一颗宝石在金黄的金座上闪亮，}
\verseline{135}{镶在颈饰或头冠上；又如匠人}
\verseline{}{把象牙嵌入黄杨木或奥里库斯的笃耨香木，}
\verseline{}{使它发光；乳白的颈项承接披散的长发，}
\verseline{}{柔软的金环在下面束住它们。}
\verseline{}{心怀壮志的族民也看见你，伊斯马鲁斯，}
\verseline{140}{射出创伤，把毒涂上箭杆；}
\verseline{}{你出身于迈俄尼亚高贵的家族，那里的人耕种}
\verseline{}{肥沃的田野，帕克托卢斯河以金砂灌溉土地。}
\verseline{}{门涅斯透斯也在场；他不久前把图尔努斯从城垣}
\verseline{}{土垒上逐退，由此赢得的荣名将他高高举起；}
\verseline{145}{还有卡皮斯，坎帕尼亚的一座城市由他得名。\footnote{守城者来自特洛伊、吕喀亚等不同支系。本卷第 145 行以卡皮斯解释坎帕尼亚城市卡普阿的名称，属于罗马诗歌常见的英雄祖源词源。}}
\verseline{}{他们在那里承受艰苦战争的交锋；}
\verseline{}{埃涅阿斯却在深夜横渡海中。}
\verseline{}{他离开厄万德，进入伊特鲁里亚营地，}
\verseline{}{面见国王，向他说明自己的姓名与族系，}
\verseline{150}{说明他寻求什么、自己带来什么，又告诉他墨泽提乌斯}
\verseline{}{聚集了怎样的兵力，图尔努斯的胸中何等暴烈；}
\verseline{}{他提醒国王凡人之事多么不足倚恃，}
\verseline{}{又把恳求杂在陈词里。塔耳孔毫不迟疑，}
\verseline{}{合并双方兵力，缔结盟约；于是吕底亚民族摆脱旧命运，}
\verseline{155}{依照神谕登上舰队，}
\verseline{}{把自己交给外邦的统帅。埃涅阿斯的船}
\verseline{}{居于最前，冲角下雕着佛里吉亚双狮，}
\verseline{}{伊达山耸立其上，最能慰藉流亡的特乌克罗斯人。\footnote{第八卷已说明伊特鲁里亚神谕：本土领袖不得统率反抗墨泽提乌斯的军队，必须由外邦人领军。埃涅阿斯加入后，这一条件满足；此处“摆脱旧命运”指解除神谕所设的拘束，吕底亚／伊特鲁里亚人因而登船出征。吕底亚称谓来自伊特鲁里亚人源出小亚细亚吕底亚的古代传统。}}
\verseline{}{伟大的埃涅阿斯坐在那里，独自反复思量}
\verseline{160}{战争多变的结局；帕拉斯紧贴在他左侧，}
\verseline{}{一会儿问星辰，一会儿问幽暗}
\verseline{}{夜里的航路，又问他在海陆经历的一切。}
\verseline{}{女神们，现在打开赫利孔，唤起歌声：}
\verseline{}{是哪一支军队从伊特鲁里亚海岸随埃涅阿斯而来，}
\verseline{165}{武装船舰，乘海航行？\footnote{舰队名录共三十艘船。诗依次点出克卢西乌姆、科萨、波普洛尼亚、伊尔瓦、皮萨、凯雷、皮耳癸、格拉维斯凯、曼图亚、贝纳库斯与明基乌斯等地，说明伊特鲁里亚联盟的来源；各船又装饰猛虎、阿波罗、半人马和特里同等船首像。}}
\verseline{}{马西库斯居首，以青铜猛虎为船首划开水面；}
\verseline{}{他麾下有一千名青年，他们离开克卢西乌姆的城墙，}
\verseline{}{也离开科萨城；他们的兵器是箭，肩上是}
\verseline{}{轻便的箭箙与致命的弓。}
\verseline{170}{凶悍的阿巴斯与他同行；整支队伍的兵甲}
\verseline{}{华美夺目，金色阿波罗在船尾闪耀。}
\verseline{}{母城波普洛尼亚给他六百名}
\verseline{}{久经战阵的青年；伊尔瓦又给三百，}
\verseline{}{那岛盛产取之不竭的哈吕柏斯金属。}
\verseline{175}{第三位是人与神的阐释者阿西拉斯；}
\verseline{}{牲畜的内脏、天上的星辰都服从他，}
\verseline{}{还有鸟语与预告将来的雷火；}
\verseline{}{他匆匆带来一千人，阵列密集，长矛森然。}
\verseline{}{命令他们服从的是源出阿尔甫斯河的皮萨，}
\verseline{180}{坐落在伊特鲁里亚土地上的城市。随后是最俊美的阿斯图尔，}
\verseline{}{阿斯图尔信赖战马，也信赖斑斓的兵甲。}
\verseline{}{又有三百人加入——众人同怀追随之心——}
\verseline{}{他们来自凯雷，来自米尼翁河的田野，}
\verseline{}{来自古老的皮耳癸与气候不善的格拉维斯凯。}
\verseline{185}{我不会略过你，战阵中最勇猛的利古里亚首领}
\verseline{}{库涅鲁斯，也不会略过只有少数人相随的库帕沃；}
\verseline{}{天鹅羽毛从他的盔顶耸起——}
\verseline{}{爱神啊，那是你的罪过，也是他父亲美丽形貌的标记。}
\verseline{}{相传库克努斯哀悼心爱的法厄同，}
\verseline{190}{在化作白杨的姐妹们枝叶与浓荫之间歌唱，}
\verseline{}{以缪斯之声安慰悲伤的爱；}
\verseline{}{在柔软的羽毛中度过霜白的晚年，}
\verseline{}{终于离开大地，以歌声追随星辰。\footnote{185—193行写库克努斯因哀悼心爱的法厄同而化为天鹅；法厄同的姐妹则化为白杨。库帕沃头盔上的天鹅羽饰指向父亲库克努斯的变形，也保存了他对法厄同的爱与歌声中的哀悼。羽饰也纪念父亲的形貌之美。192行 canentem 来自 caneo（呈白色），修饰 senectam（晚年），以霜白联系天鹅羽色。Amor 在诗中被直接呼告；Conte 2019 校勘记将188行 vestrum 解作单数 tuum，呼告对象仍是单数的爱神。}}
\verseline{}{他的儿子带着同龄的队伍，}
\verseline{195}{以桨推动巨大的“半人马”号；船首像}
\verseline{}{俯压水面，以巨岩威胁波涛，}
\verseline{}{高高耸立，长长的龙骨犁开深海。}
\verseline{}{奥克努斯也从故乡海岸召起队伍，}
\verseline{}{他是女先知曼托与伊特鲁里亚河神之子；}
\verseline{200}{曼图亚，是他给你城墙，也以母亲为你命名，}
\verseline{}{曼图亚，先祖显赫，却并非所有人同出一族：}
\verseline{}{那里有三个族系，每族下分四个部族，}
\verseline{}{城本身为诸部族之首，力量来自伊特鲁里亚血统。}
\verseline{}{墨泽提乌斯也从那里为自己招来五百名武装的敌人，}
\verseline{205}{明基乌斯头戴灰绿的芦苇，由父亲贝纳库斯所生，}
\verseline{}{驾着迎敌的松木船，把他们领向海面。\footnote{奥克努斯是女先知曼托之子，诗中说曼图亚由母名得名，又把城中居民分成三个族系、每族四个部族。贝纳库斯是今加尔达湖，明基乌斯河从湖中流出，所以诗中称河为湖之子。}}
\verseline{}{沉重的奥勒斯忒斯驶来，一百支桨扬起}
\verseline{}{击打波涛，翻转的水面白沫四起。}
\verseline{}{巨大的特里同载着他，以海螺}
\verseline{210}{惊吓青蓝的海峡；船首像游动时，直到两胁}
\verseline{}{仍显出毛茸茸的人形，腹部以下却收成海兽鱼身，}
\verseline{}{泡沫在半兽的胸膛下低鸣。}
\verseline{}{这么多精选的首领乘三十艘船而来，}
\verseline{}{支援特洛伊，以青铜划开盐海的平原。}
\verseline{215}{此时白昼已从天上退去，温柔的福柏}
\verseline{}{驾着夜游车，行至奥林匹斯中天；}
\verseline{}{忧虑不让埃涅阿斯的四肢得到休息，}
\verseline{}{亲自坐着掌舵，也料理风帆。\footnote{Conte 2019 在第 218 行后、第 219 行前排出一处缺文标记，采 Dionisotti 的判断。该版校勘记指出，第 219 行的代词 illi（指埃涅阿斯的座船）以及 suarum ... comitum（指同为旧特洛伊舰船的宁芙同伴）缺少前文明确先行词，因此推测此处曾失去一行。}}
\verseline{}{就在航程中途，看哪，一群旧日的同伴}
\verseline{220}{迎面而来：大母神库柏柏曾命令}
\verseline{}{她们从舰船变为宁芙，获得海的神性；}
\verseline{}{她们并排游来，一齐划开波浪，}
\verseline{}{数目正如从前立在海岸的青铜船首。}
\verseline{}{她们远远认出国王，以环舞围绕他；}
\verseline{225}{其中最善言辞的库摩多刻亚}
\verseline{}{从后方跟随，右手扶住船尾，背脊}
\verseline{}{露出水面，左手在寂静的波下划动，}
\verseline{}{随后对毫不知情的他开口：“神的后裔}
\verseline{}{埃涅阿斯，你醒着吗？醒来，把帆索放开。}
\verseline{230}{我们是从伊达山神圣峰顶伐下的松木，}
\verseline{}{从前是你的舰队，如今是海上的宁芙。背信的}
\verseline{}{鲁图利人以剑与火催迫我们时，}
\verseline{}{我们不情愿地挣断你的缆绳，穿过海面寻找你。}
\verseline{}{众神之母怜悯我们，重塑这副形体，}
\verseline{235}{又使我们成为女神，在波下度过永世。}
\verseline{}{而少年阿斯卡尼乌斯仍被困在城墙和壕沟内，}
\verseline{}{处在兵刃中央，面对因战神而森然的拉丁人。}
\verseline{}{阿卡狄亚骑兵已与强悍的伊特鲁里亚人会合，}
\verseline{}{占据指定的位置；图尔努斯决意把骑兵队}
\verseline{240}{横挡在他们面前，不让他们与营地会合。}
\verseline{}{起来吧；曙光一来，你要首先下令}
\verseline{}{召集同伴武装，拿起火之主亲自赐予的}
\verseline{}{无敌盾牌——黄金环绕它的边缘。}
\verseline{}{若你不以为我的话虚妄，明日的晨光}
\verseline{245}{将看见鲁图利人尸体堆成巨丘。”\footnote{舰船神化承接第九卷：伊达山大母神圣木所造之船在鲁图利人纵火时挣断缆绳，化为海中宁芙。库摩多刻亚（Cymodocea）是其中最善言辞者；第一卷的海宁芙库摩托厄（Cymothoe）与她词形相近，实为另一位神女。她右手托船，左手在水下划动，在游行中向埃涅阿斯传话。}}
\verseline{}{她说完，离开时以右手推动高高的船尾，}
\verseline{}{深谙其法；船便越过水面飞驰，}
\verseline{}{比标枪还快，比追上疾风的箭更迅急。}
\verseline{}{其余宁芙也加快航程。安喀塞斯之子}
\verseline{250}{这位特洛伊人惊奇不已，却因吉兆振作心神。}
\verseline{}{随后他仰望穹苍，作了简短祈祷：}
\verseline{}{“仁慈的伊达山众神之母，你爱丁底马山，}
\verseline{}{爱头戴塔冠的城市，也爱缰绳下成对的雄狮；}
\verseline{}{如今请先引我进入战斗，让这预兆依礼到来，}
\verseline{255}{以顺遂的脚步临助佛里吉亚人，女神。”\footnote{“伊达山众神之母”即大母神库柏勒；220行的库柏柏（Cybebe）是同神的另一名称，对应 Cybele／Cybebe。丁底马山、塔冠及成对雄狮都是她的传统标志；253行的城楼城市也与塔冠形象相连。}}
\verseline{}{他只说了这些。与此同时，轮转而来的白昼正疾驰，}
\verseline{}{天光已经大亮，也已将黑夜驱散。}
\verseline{}{他先命同伴跟随军旗，}
\verseline{}{让战心与兵器一同就绪，准备交战。}
\verseline{260}{他站在高高的船尾，已看见}
\verseline{}{特乌克罗斯人和自己的营地；随后以左手}
\verseline{}{举起燃烧的盾牌。达耳达尼亚人从城墙上}
\verseline{}{发出直达星辰的呼喊，新生的希望激起怒火，}
\verseline{}{他们手掷兵刃；犹如黑云下}
\verseline{265}{斯特律蒙河的鹤群发出信号，鸣叫着}
\verseline{}{游过长空，以欣然的呼声逃避南风。}
\verseline{}{鲁图利人的王与奥索尼亚诸将却感到惊异，}
\verseline{}{直到他们回头望见舰船一齐转向岸边，}
\verseline{}{整片海仿佛随舰队滑向陆地。}
\verseline{270}{埃涅阿斯的盔顶燃烧，阴森的火焰从头顶泻下，}
\verseline{}{金色盾脐吐出巨大的火；}
\verseline{}{正如澄澈夜空中血红的彗星}
\verseline{}{发出哀凶的红光，又如天狼星升起，}
\verseline{}{为患病的凡人带来干渴与瘟疫，}
\verseline{275}{以不祥的光使天空悲暗。\footnote{神盾与盔顶的火光先被比作血红彗星，再比作带来干渴、疾病的天狼星。它对守军是埃涅阿斯归来的信号，对敌我双方却都保留凶兆和灾疫的色彩。}}
\verseline{}{但勇敢的图尔努斯没有失去信心，}
\verseline{}{仍要抢占海岸，把来者赶离陆地。}
\verseline{}{他反而以言语振作众心，反而斥责他们：\footnote{Conte 2019 将第 278 行以方括号排印；该行重现第九卷第 127 行，该版校勘记指出若干主要传本缺失。}}
\verseline{}{“你们祈求的时刻已经来到——以右手将他们粉碎；}
\verseline{280}{战神就在手中，勇士们！现在，每个人都要记起}
\verseline{}{自己的妻子和家宅，记起先辈的伟绩}
\verseline{}{与声名。让我们主动到水边迎击，}
\verseline{}{趁他们慌乱登陆、第一步尚未站稳。}
\verseline{}{命运帮助敢于行动的人。”\footnote{Audentis Fortuna iuvat处在图尔努斯部署反登陆战的军令中，意为“命运帮助敢于行动的人”。这里的Fortuna指战场时运，图尔努斯借此催促部下抢占滩头。马克罗比乌斯（6.1.62）又把这句与恩尼乌斯《编年史》第七卷Fortibus est fortuna viris data（“命运赐给勇者”）并列。}}
\verseline{285}{他说完，暗自思量该带哪些人迎战，}
\verseline{}{又能把受围的城墙托付给谁。}
\verseline{}{与此同时，埃涅阿斯架起桥板，让同伴从高船下来。}
\verseline{}{许多人看准退潮，}
\verseline{}{趁海水渐弱，纵身跃入浅滩；}
\verseline{290}{另一些人顺着船桨下水。塔耳孔察看海岸，}
\verseline{}{寻找浅滩不作喘鸣、碎浪也不回响之处，}
\verseline{}{那里海面平顺，随涨潮无碍地滑向岸边；}
\verseline{}{他忽然调转船首，向同伴祈请：}
\verseline{}{“精锐的勇士们，现在压上强健的船桨，}
\verseline{295}{抬起船，送它向前；让冲角劈开这敌对的}
\verseline{}{土地，让龙骨亲自犁入土中。}
\verseline{}{即使船尾撞碎在这样的泊地，我也在所不惜，}
\verseline{}{只要一举夺得陆地。”塔耳孔说完这番话，}
\verseline{}{同伴们一齐压上桨柄，}
\verseline{300}{把白沫四溅的船驶向拉丁田野，}
\verseline{}{直到冲角抵住干地，所有龙骨}
\verseline{}{都安然坐上海岸——惟独你的船不是如此，塔耳孔。}
\verseline{}{它撞入浅滩，悬在一道高低不平的沙脊上，}
\verseline{}{摇摆着支撑许久，承受波浪反复冲击，}
\verseline{305}{终于裂开，把船上的人抛进水中；}
\verseline{}{断桨与漂浮的座板缠住他们，}
\verseline{}{退去的波浪又同时把他们拖回海里。\footnote{塔耳孔选择潮水平稳抵岸、没有碎浪回声的地段，命船只直接以龙骨冲上陆地，宁可以船尾破裂换取登岸。其他船安全搁浅，他的船却架在不平沙脊上，终于裂开。304行 fluctus 可按主格单数理解为波浪冲击船只，塞尔维乌斯（10.304）采用这一句法；也可按宾格复数理解为船只抗击波浪。退潮、浅滩、桥板、断桨和漂浮座板共同写出登陆与船毁后的困境。}}
\verseline{}{但迟延不能使图尔努斯怠慢；他勇猛地}
\verseline{}{把整个战阵推向特乌克罗斯人，在海岸上列阵相对。}
\verseline{310}{军号响起。埃涅阿斯率先冲入乡野的队列，}
\verseline{}{为战斗立下预兆；他击溃拉丁人，}
\verseline{}{杀死忒戎——那身躯最高大的勇士主动}
\verseline{}{迎向埃涅阿斯；埃涅阿斯以剑穿过铜片缝成的甲胄，}
\verseline{}{穿过金光森然的战衣，深深刺入他暴露的胁腹。}
\verseline{315}{接着他击中利卡斯；母亲死去时，人们曾剖腹取出他，}
\verseline{}{把他献给你，福玻斯；幼小的他得以逃过刀锋，}
\verseline{}{如今又有什么用？不远处，基修斯与}
\verseline{}{巨大的古阿斯挥舞棍棒，扫倒成排敌军，}
\verseline{}{却也被他送入死亡；赫尔库勒斯的兵器、}
\verseline{320}{他们强健的双手与父亲墨兰普斯都未能相救——}
\verseline{}{只要大地还给阿尔喀得斯严酷的劳役，}
\verseline{}{墨兰普斯便一直伴随他。看哪，法罗斯还在抛出无力的狂言，}
\verseline{}{埃涅阿斯便扭身投枪，把矛钉进他呼喊的口中。}
\verseline{}{不幸的库冬，你追逐脸颊初生}
\verseline{325}{金黄绒毛的克吕提乌斯——你新得的欢愉——}
\verseline{}{原本也会被达耳达尼亚人的右手击倒，}
\verseline{}{可怜地躺在那里，不再忧虑你从来眷爱的青年男子，}
\verseline{}{若不是一队密集的兄弟迎面赶来——}
\verseline{}{他们是福尔库斯的七个儿子，投来七支兵刃；}
\verseline{330}{一部分无效地从头盔和盾牌弹回，}
\verseline{}{另一部分擦近身体，却被仁慈的维纳斯拨开。\footnote{本卷第 324—331 行写库冬追逐脸颊初生金黄绒毛的青年克吕提乌斯；原诗称后者为“新近的欢愉”，并说库冬一向爱恋青年男子。这段爱欲关系被置于史诗战斗名录之中。}}
\verseline{}{埃涅阿斯对忠诚的阿卡忒斯说：}
\verseline{}{“递给我兵器；这些曾在伊利乌姆平原上}
\verseline{}{刺进希腊人身体的枪，我的右手不会有一支}
\verseline{335}{白白投向鲁图利人。”他抓起一柄巨矛}
\verseline{}{投出；长矛飞去，贯穿迈翁盾上的青铜，}
\verseline{}{又撕裂胸甲，穿入胸膛。}
\verseline{}{迈翁的兄弟阿尔卡诺耳赶来，以右手}
\verseline{}{扶住倒下的人；长矛穿透阿尔卡诺耳的手臂，}
\verseline{340}{仍带着鲜血向前飞去，保持原来的去势；}
\verseline{}{右手垂死地悬在肩下，只由筋腱牵连。}
\verseline{}{努米托耳从兄弟身上拔出长矛，}
\verseline{}{又以它袭击埃涅阿斯；却未获准}
\verseline{}{把矛刺进对手，只擦伤伟大的阿卡忒斯的大腿。}
\verseline{345}{这时，来自库瑞斯、正信赖青春体魄的克劳苏斯}
\verseline{}{赶到，从远处以坚硬长矛击中德律俄普斯，}
\verseline{}{枪尖重重抵在下颏，刺穿喉咙，}
\verseline{}{一并夺去他说话的声音与生命；}
\verseline{}{他额头撞地，口中吐出浓稠的血。}
\verseline{350}{克劳苏斯又以不同的死法杀死三个}
\verseline{}{出自玻瑞阿斯最高血统的色雷斯人，}
\verseline{}{以及父亲伊达斯与祖国伊斯马鲁斯送来的三人。}
\verseline{}{哈莱苏斯和奥伦卡军队奔来，}
\verseline{}{涅普顿之子、以战马闻名的墨萨波斯也迎上前。}
\verseline{355}{这一方、那一方轮番竭力驱逐对手；}
\verseline{}{双方就在奥索尼亚的门槛上争战，仿佛广天中相冲的风}
\verseline{}{以相等的意气和力量掀起交锋；}
\verseline{}{它们彼此不退，云不退，海也不退，}
\verseline{}{战斗长久未决，万物相持对抗：}
\verseline{360}{特洛伊战阵与拉丁战阵也这样}
\verseline{}{冲撞；脚抵住脚，密集的人抵住人。}
\verseline{}{战场另一处，一道急流把滚石冲向远方，}
\verseline{}{把树丛连根卷离堤岸；}
\verseline{}{帕拉斯看见不惯步战的阿卡狄亚人}
\verseline{365}{背向追赶的拉丁人——}
\verseline{}{崎岖的地势迫使他们舍弃战马——}
\verseline{}{危急中只剩下这一条路：他时而恳求，}
\verseline{}{时而以严厉的话激起勇气：}
\verseline{}{“同伴们，你们逃向哪里？凭你们与自己的英勇事迹，}
\verseline{370}{凭厄万德王的姓名与他赢得的战争，}
\verseline{}{也凭我如今要与父辈声名相竞的希望，}
\verseline{}{不要信赖双脚。必须以剑破开}
\verseline{}{穿过敌军的道路。敌群最密集地逼来的地方，}
\verseline{}{正是高贵的祖国召唤你们与领袖帕拉斯之处。}
\verseline{375}{没有神力压迫我们；我们这些凡人受凡人仇敌逼迫，}
\verseline{}{双方有同样多的生命、同样多的手。}
\verseline{}{看哪，海以巨大的阻障封住我们，}
\verseline{}{陆上已没有逃路：我们要下海，还是返回特洛伊？”}
\verseline{}{他说完，从中间突入密集的敌阵。}
\verseline{380}{最先迎上他的是拉古斯，}
\verseline{}{被不利的命运引来；他正用巨大力量撬动岩石，}
\verseline{}{帕拉斯便扭身投枪，把他钉在脊椎}
\verseline{}{从中央隔开两侧肋骨之处，又拔回}
\verseline{}{卡在骨中的长矛。希斯波未能趁势压倒他，}
\verseline{385}{虽说他抱着这份指望；帕拉斯抢先迎住来者，}
\verseline{}{趁他因同伴的残酷死亡狂怒、毫无防备，}
\verseline{}{把剑埋进他肿胀的肺中。}
\verseline{}{随后他攻向斯忒纽斯，又攻向罗厄图斯古老家族的}
\verseline{}{安刻摩卢斯——他竟敢玷污继母的婚床。}
\verseline{390}{你们这对孪生子也倒在鲁图利人的田野，}
\verseline{}{达乌库斯之子拉里斯与提姆布鲁斯；你们如此相似，}
\verseline{}{亲人也无法分辨，父母曾喜爱那甜美的错认；}
\verseline{}{如今帕拉斯却给你们严酷的区别：}
\verseline{}{提姆布鲁斯，厄万德之子的剑斩去了你的头；}
\verseline{395}{拉里斯，你被斩断的右手仍寻找主人，}
\verseline{}{半死的手指跳动着，要重新抓住剑柄。}
\verseline{}{阿卡狄亚人因他的告诫燃起勇气，又看见}
\verseline{}{这位英雄的卓越战绩；悲痛与羞耻交织，武装他们冲入敌阵。}
\verseline{}{随后帕拉斯投枪，贯穿驾双马车逃走的瑞图斯；}
\verseline{400}{这一阻隔只给伊卢斯片刻的迟延，}
\verseline{}{因为他原从远处把强矛投向伊卢斯，}
\verseline{}{瑞图斯却挡在中间——最英勇的透特拉斯啊，}
\verseline{}{他正在逃避你和你的兄弟提瑞斯——从车上滚落，}
\verseline{}{半死的脚跟击打鲁图利人的土地。}
\verseline{405}{犹如盛夏，期望已久的风一旦升起，}
\verseline{}{牧人便把零散火种送入林间；}
\verseline{}{火焰突然咬住中间的树木，}
\verseline{}{连成火神森严的一线，延伸过辽阔原野，}
\verseline{}{牧人坐在那里，以胜利者之姿俯视欢腾的火：}
\verseline{410}{你的同伴也汇聚全部勇力，}
\verseline{}{前来援助你，帕拉斯。但战场上刚烈的哈莱苏斯}
\verseline{}{迎向他们，把自己收拢在盾甲之后。}
\verseline{}{他杀死拉东、费瑞斯和得摩多库斯；}
\verseline{}{斯特里蒙举起右手，直取他的咽喉，}
\verseline{415}{他却以闪亮的剑将手斩落；又以岩石击中托阿斯的面孔，}
\verseline{}{击得骨片与带血的脑髓一同迸散。}
\verseline{}{哈莱苏斯的父亲预知命运，曾把他藏进森林；}
\verseline{}{当老人死去，合上预见死亡的眼睛，}
\verseline{}{帕耳卡三女神便把手伸向哈莱苏斯，}
\verseline{420}{将他献给厄万德之子的兵刃。帕拉斯进攻以前这样祈祷：\footnote{哈莱苏斯的父亲预知儿子命运，生前把他藏在森林；父亲死后，帕耳卡三女神才把哈莱苏斯交给战争，并“献定”给厄万德之子的兵器。祭献语汇将死亡同预知的命运相连。418行 canentia lumina 可承上作“预言命运的眼睛”；塞尔维乌斯（10.418）也列临死瞳孔转白的读法。}}
\verseline{}{“父亲台伯河啊，请把运势与道路赐给我正要掷出的铁枪，}
\verseline{}{让它穿过哈莱苏斯坚硬的胸膛。}
\verseline{}{你的橡树将悬挂这个人的兵甲与战利品。”}
\verseline{}{河神听见了；哈莱苏斯正遮护伊马翁，}
\verseline{425}{不幸把无甲的胸膛交给阿卡狄亚人的枪。\footnote{帕拉斯向人格化的台伯河神祈求长矛获得“运势与道路”，并许诺在杀死哈莱苏斯后，把从他身上夺下的兵甲悬在河神的橡树上。他预先把战果许给河神，作为还愿之物；这些兵甲首先属于祭献。}}
\verseline{}{劳苏斯却没有因这名勇士的死亡而惊慌；}
\verseline{}{他是战争的重要支柱，不许队伍溃散，首先杀死}
\verseline{}{迎面的阿巴斯——战斗的死结与阻碍。}
\verseline{}{阿卡狄亚的子弟倒下，伊特鲁里亚人倒下，}
\verseline{430}{你们也倒下了，特乌克罗斯人——希腊人未能毁灭的躯体。}
\verseline{}{两军相遇，领袖与力量彼此相等。}
\verseline{}{后列挤紧战阵，人群拥塞，}
\verseline{}{挤得手臂与兵刃全都动弹不得。帕拉斯从这一边进逼，}
\verseline{}{劳苏斯从另一边抵抗；二人年岁相差无几，}
\verseline{435}{容貌同样出众，命运却不许他们}
\verseline{}{返回故乡。伟大奥林匹斯的统治者}
\verseline{}{也不准他们彼此交锋；}
\verseline{}{不久，各自的命运将在更强大的敌人手下等候他们。\footnote{帕拉斯与劳苏斯年岁相近、容貌出众，都不能还乡。朱庇特阻止他们彼此交战；后来帕拉斯死于图尔努斯之手，劳苏斯死于埃涅阿斯之手。}}
\verseline{}{与此同时，图尔努斯神圣的姊妹劝他接替劳苏斯；\footnote{原文只称图尔努斯的“神圣姊妹”；这一称号通常指后来正式出场的河泉女神尤图耳娜，也延迟了她的命名。}}
\verseline{440}{他驾着飞驰的战车穿过队伍中央。}
\verseline{}{一看见同伴，他便说：“该停下战斗了；}
\verseline{}{只有我来迎战帕拉斯，帕拉斯只应归我。}
\verseline{}{但愿他的父亲亲自在这里观看。”}
\verseline{}{他说完，同伴遵令退出平野。}
\verseline{445}{鲁图利人一退开，青年便对这傲慢的命令}
\verseline{}{惊愕不已，怔怔望着图尔努斯；目光沿他庞大的}
\verseline{}{身躯游移，从远处以凶狠的眼神把他上下打量，}
\verseline{}{便用这番话回敬暴君的话：}
\verseline{}{“我今日要么夺得敌帅的兵甲，赢得最高赞誉，}
\verseline{450}{要么以显赫之死赢得声名；父亲对两种命运都能坦然承受。\footnote{马克罗比乌斯（6.1.56）把帕拉斯这句“要么夺得战利品，要么获显赫之死”与阿克齐乌斯失传悲剧《兵器审判》（Armorum iudicium）并列：从强者手中取得战利品固然光荣，若败于这样的强者也不算耻辱。}}
\verseline{}{收起威胁吧。”他说完，走向平野中央；}
\verseline{}{冰冷的血在阿卡狄亚人胸中凝结。}
\verseline{}{图尔努斯从双马车跃下，准备徒步}
\verseline{}{近战；犹如狮子从高处望见}
\verseline{455}{远处田野上的公牛站定、预备交战，}
\verseline{}{便飞扑而来：图尔努斯逼近的形象正是如此。}
\verseline{}{帕拉斯相信他已来到投枪可及之处，}
\verseline{}{便抢先向前，希望某种机运能帮助他}
\verseline{}{以不相等的力量行此壮举；他向高天祷告：}
\verseline{460}{“阿尔喀得斯啊，凭我父亲与你的宾主之谊和宴席，}
\verseline{}{求你临助我的宏举。}
\verseline{}{让图尔努斯在半死之际看见我夺取他染血的兵甲，}
\verseline{}{让他将死的眼睛承认我是胜者。”\footnote{帕拉斯凭厄万德与赫尔库勒斯的宾主关系和共同宴席向阿尔喀得斯祈求援助。第八卷中厄万德家族以祭礼纪念赫尔库勒斯杀死卡库斯；此处赫尔库勒斯听见祷告却只能流下“徒劳的泪”：旧日宾主关系仍在，神却无力救下帕拉斯。}}
\verseline{}{阿尔喀得斯听见青年的祈求，把一声深长的叹息}
\verseline{465}{压在心底，流下徒劳的眼泪。}
\verseline{}{于是众神之父以亲切的话对儿子说：}
\verseline{}{“每个人的死期都已规定；众生的生命}
\verseline{}{短暂而不可挽回，但以功业延长声名，}
\verseline{}{这才是勇武的工作。特洛伊高墙之下}
\verseline{470}{曾有多少神的儿子倒下；我的儿子}
\verseline{}{萨尔佩冬不也一同死去？图尔努斯也正被}
\verseline{}{自己的命运召唤，已走到所赐年岁的尽头。”\footnote{Stat sua cuique dies 指每个人都有既定死期；凡人只能借功业延长 fama（“声名”）。朱庇特以萨尔佩冬说明神之子同样受制于死亡。他在帕拉斯与图尔努斯决斗前说出这番话，并预告图尔努斯的终点逼近；帕拉斯随后丧命，厄万德也将承受丧子之痛。}}
\verseline{}{他说完，又把目光转回鲁图利人的田野。}
\verseline{}{帕拉斯用尽力气掷出长矛，}
\verseline{475}{又从剑鞘中拔出闪亮的剑。}
\verseline{}{长矛飞去，击在肩甲隆起的最高处，}
\verseline{}{勉强从盾缘开出道路，}
\verseline{}{终于擦伤图尔努斯庞大的身躯。}
\verseline{}{图尔努斯把尖铁装在前端的橡木长矛}
\verseline{480}{长久瞄准帕拉斯，掷出时说道：}
\verseline{}{“现在看看，我的兵器是否穿得更深。”}
\verseline{}{他说完；矛尖震动着一击，竟贯穿盾牌中央——}
\verseline{}{多少层铁，多少层青铜，}
\verseline{}{多少层裹覆其上的牛皮——}
\verseline{485}{它又突破胸甲的阻拦，穿透宽阔的胸膛。}
\verseline{}{青年徒然从伤口拔取滚烫的枪；}
\verseline{}{鲜血与生命沿同一条道路流出。}
\verseline{}{他倒在伤口上，兵甲轰然作响，}
\verseline{}{临死以染血的口触到敌人的土地。}
\verseline{490}{图尔努斯站在他上方，}
\verseline{}{说道：“阿卡狄亚人，记住我的话，带给}
\verseline{}{厄万德：我把他应得的帕拉斯还给他。}
\verseline{}{坟冢应有的一切荣誉、下葬的一切慰藉，}
\verseline{}{我都慷慨赐下。他与埃涅阿斯结下的}
\verseline{495}{宾主之谊不会只让他付出小小代价。”说完，他以左脚踩住死者，}
\verseline{}{从尸体上扯下那沉重的剑带；带面刻着}
\verseline{}{神律所禁之事：同一新婚之夜，一队青年}
\verseline{}{遭到污秽的屠杀，婚床染血；}
\verseline{}{欧律图斯之子克洛努斯曾用大量黄金雕成这幅图景。}
\verseline{500}{图尔努斯如今为这件战利品欢欣，为占有它而喜悦。\footnote{剑带斜挎肩胸，用来悬佩武器；第十二卷结尾，同一剑带会在图尔努斯肩头重现。带上图案通常被解释为达那伊得斯在新婚夜杀死丈夫：诗写青年被杀、婚床染血，并称带上的这一画面为 nefas（“神律所禁”）。胜者依战场惯例可剥取死者兵甲；图尔努斯却踩住帕拉斯尸体，夺下并披上剑带，诗人随即预言他终将憎恨这件战利品与这一天。}}
\verseline{}{凡人的心不知命运，不知未来的遭际，}
\verseline{}{一旦被顺境托起，也不知谨守分寸！}
\verseline{}{总有一天，图尔努斯愿以重价换取}
\verseline{}{帕拉斯不受触碰，也将憎恨这件战利品与这一天。}
\verseline{505}{同伴们却在长久的哀号和泪水中聚拢，}
\verseline{}{把帕拉斯放上盾牌，抬着他回去。}
\verseline{}{啊，你将带给父亲何等巨大的悲痛与荣光！}
\verseline{}{这是战争把你交出的第一天，也是把你夺走的一天，}
\verseline{}{而你仍留下鲁图利人尸体堆成的巨丘。\footnote{同伴以盾牌充作临时尸床，把帕拉斯抬回。诗人直接呼告青年：出战第一天即死亡之日；他给父亲带来荣光，也带来深重的悲痛。“留下敌尸”的战场荣誉与父亲的丧子之痛并置。}}
\verseline{510}{来到埃涅阿斯耳边的已不是如此巨祸的传闻，}
\verseline{}{而是确切的报信者：他的人处在细若一线的生死之间，}
\verseline{}{该去援救正在后退的特乌克罗斯人了。}
\verseline{}{他以剑割倒离自己最近的一切，在宽阔的敌阵中}
\verseline{}{燃烧着以铁开路，寻找你，图尔努斯，}
\verseline{515}{因新近杀戮而骄傲的你。帕拉斯、厄万德，}
\verseline{}{一切都在他眼前：他初作宾客时走近的餐桌，}
\verseline{}{还有彼此相握的右手。他活捉四个}
\verseline{}{苏尔摩之子，又活捉乌芬斯养大的四人，}
\verseline{}{要把这些活人献祭给亡魂，}
\verseline{520}{以俘虏的血浇淋火葬堆的火焰。\footnote{埃涅阿斯活捉苏尔摩的四个儿子和乌芬斯养大的四名青年，准备在帕拉斯葬礼上杀死，以俘虏的血浇洒葬火，作为献给亡魂的祭品。可参《伊利亚特》第二十三卷，阿喀琉斯为帕特罗克洛斯杀死十二名特洛伊青年。}}
\verseline{}{接着，他从远处把敌意的长矛瞄向玛古斯。}
\verseline{}{玛古斯机敏地伏下，枪身颤动着从上方飞过；}
\verseline{}{他抱住埃涅阿斯的膝盖，像恳求者一样开口：}
\verseline{}{“凭你父亲的亡魂和正在长大的尤卢斯的希望，}
\verseline{525}{求你为我的儿子与父亲保全这条生命。}
\verseline{}{我有高大的宅邸，那里深埋着数塔兰同的}
\verseline{}{雕花白银；我也有大量黄金，}
\verseline{}{有已经锻造的，也有未经锻造的。特乌克罗斯人的胜利}
\verseline{}{不会因此改变，我一条命不会造成如此大的差别。”}
\verseline{530}{他说完。埃涅阿斯这样回答：}
\verseline{}{“你所说的许多塔兰同的金银，}
\verseline{}{留给你的儿子吧。图尔努斯杀死帕拉斯时，}
\verseline{}{早已断绝这种战场上的赎命交易。}
\verseline{}{我父亲安喀塞斯的亡魂这样认为，尤卢斯也这样认为。”\footnote{玛古斯以儿孙、金银与战争胜负求饶，埃涅阿斯回答图尔努斯在杀死帕拉斯时已经断绝“战场上的赎命交易”，随后杀死恳求者。他还说安喀塞斯的亡魂和尤卢斯都不允许自己饶命，把杀死恳求者说成替死者与家人复仇。}}
\verseline{535}{他说完，以左手抓住头盔，把恳求者的颈项}
\verseline{}{向后扳去，将剑直刺到柄端。}
\verseline{}{不远处是海蒙尼得斯，福玻斯与特里维亚的祭司，}
\verseline{}{神圣的额带和圣带环绕他的两鬓，}
\verseline{}{白色衣袍与洁白的圣饰使他通体闪亮。}
\verseline{540}{埃涅阿斯迎上他，将他逐过平野；祭司倒下后，他立在上方，}
\verseline{}{把他如祭牲般宰杀，以巨大的阴影覆盖；塞雷斯图斯}
\verseline{}{把精选的兵甲背回，作为给你的战胜纪念，格拉狄乌斯王。}
\verseline{}{武尔坎努斯之子凯库卢斯与}
\verseline{}{从玛耳西人山中来的翁布罗重整阵线；}
\verseline{545}{达耳达尼亚的英雄却在对面狂怒。他以剑斩下安克苏尔的左手，}
\verseline{}{连盾牌的整个圆面也一并击落在地；}
\verseline{}{安克苏尔曾说过什么豪言，以为力量会支持这话，}
\verseline{}{也许他的心已升向天际，}
\verseline{}{许诺自己白发与漫长的岁月。}
\verseline{550}{塔耳奎图斯身披闪耀的兵甲，欢腾地}
\verseline{}{迎向狂怒的埃涅阿斯；林中仙女德律俄珀}
\verseline{}{为林居的法乌努斯生下了他。埃涅阿斯向后扬起长矛，}
\verseline{}{一枪将他的胸甲与沉重的盾牌钉在一起；}
\verseline{}{他还在徒然祈求，准备说许多话，}
\verseline{555}{埃涅阿斯便击落他的头，让温热的躯干}
\verseline{}{滚在地上，又怀着敌意说道：}
\verseline{}{“就躺在那里吧，可怕的勇士。你那最好的母亲}
\verseline{}{不会把你葬入土中，让你的肢体安息在祖墓；}
\verseline{}{你将留给凶禽啄食，或沉入漩涡，}
\verseline{560}{随波漂去，让饥饿的鱼舔舐伤口。”}
\verseline{}{随即他追赶安泰乌斯与卢卡斯——图尔努斯的前列——}
\verseline{}{又追赶强悍的努玛和金发卡墨耳斯，}
\verseline{}{他由心怀壮志的沃尔肯斯所生，是奥索尼亚人中}
\verseline{}{土地最富有的人，统治寂静的阿密克莱。}
\verseline{565}{犹如埃该翁——传说他有一百条臂，}
\verseline{}{一百只手，五十张口和胸膛}
\verseline{}{喷出火焰；他对抗朱庇特的雷霆时，}
\verseline{}{让同样多的盾牌一齐轰响，拔出同样多的剑——}
\verseline{}{埃涅阿斯也这样以胜者之姿在整片平野施展狂暴，}
\verseline{570}{一旦剑锋温热便不止息。看哪，他又朝尼法埃乌斯的}
\verseline{}{四马车直冲，迎向战马的胸膛。}
\verseline{}{马从远处看见他阔步走来，可怖地咆哮，}
\verseline{}{因恐惧转身，向后冲去，}
\verseline{}{把主人抛下，让车急驰向海岸。}
\verseline{575}{与此同时，卢卡古斯驾白色双马车冲入人群，}
\verseline{}{兄弟利革尔与他同车；利革尔以缰绳}
\verseline{}{引导战马，勇猛的卢卡古斯挥舞出鞘之剑。}
\verseline{}{埃涅阿斯不能容忍他们以如此炽烈的狂态进逼；}
\verseline{}{他冲上前，巨大的长矛迎面出现。}
\verseline{580}{利革尔对他说：}
\verseline{}{“你看见的不是狄俄墨得斯的战马、阿喀琉斯的战车，}
\verseline{}{也不是佛里吉亚平原；你的战争与生命}
\verseline{}{将在这片土地上结束。”疯狂的利革尔}
\verseline{}{让这些话飞过平野；特洛伊英雄}
\verseline{585}{没有准备话语回答，只把长枪掷向敌人。}
\verseline{}{卢卡古斯俯身向前，用剑身催促双马，}
\verseline{}{伸出左脚，准备迎战；}
\verseline{}{长枪从闪亮盾牌的下缘进入，}
\verseline{}{继而刺穿左侧腹股沟；}
\verseline{590}{他被撞出战车，垂死着滚在田野。}
\verseline{}{虔敬的埃涅阿斯向他投去辛辣的话：\footnote{在复仇狂暴、活人祭献和数次拒绝求饶之后，叙述者仍使用固定称号 pius Aeneas。这个“虔敬的埃涅阿斯”紧接着他的讥刺与杀戮出现，使固定称号在血腥场景中显得格外尖锐。}}
\verseline{}{“卢卡古斯，战马没有以怯懦逃跑来出卖你的战车，}
\verseline{}{敌人的空幻影子也没有使它们转身；}
\verseline{}{你亲自从车轮旁跳下，抛弃了车轭。”}
\verseline{595}{他说着攫住双马；卢卡古斯的兄弟也从同一辆车上}
\verseline{}{跌落，不幸地伸出无力的双手：}
\verseline{}{“凭你自己，凭生下你这般英雄的父母，}
\verseline{}{特洛伊勇士，留下我的生命，怜悯这恳求者吧。”}
\verseline{}{他还在说更多话，埃涅阿斯回答：“刚才你说的}
\verseline{600}{话可完全不同。死吧，兄弟，不要离开你的兄弟。”}
\verseline{}{随后以剑锋打开胸膛中灵魂的藏处。}
\verseline{}{达耳达尼亚统帅在平野制造这样的死亡，}
\verseline{}{像奔流的洪水、像黑色旋风般狂怒。}
\verseline{}{终于，他们冲出营垒，离开营地——}
\verseline{605}{少年阿斯卡尼乌斯和那群徒然被围的青年。}
\verseline{}{与此同时，朱庇特主动对朱诺说：}
\verseline{}{“我的姊妹，也是我最亲爱的妻子，}
\verseline{}{正如你所料，你的判断没有错，}
\verseline{}{支撑特洛伊人力量的是维纳斯；勇士在战斗中强健的右手、}
\verseline{610}{刚烈而耐险的心，看来全都无足轻重。”}
\verseline{}{朱诺顺从地答道：“最美好的丈夫啊，}
\verseline{}{我已痛苦，又畏惧你严厉的话，你为何还使我忧惧？}
\verseline{}{如果我在你的爱中仍有从前曾有、也理应有的}
\verseline{}{力量，全能的神，你便不会拒绝我这一件事：}
\verseline{615}{让我把图尔努斯带离战场，}
\verseline{}{完整地保全给他父亲道努斯。}
\verseline{}{如今就让他死，让他以尽责之血向特乌克罗斯人偿罪。}
\verseline{}{可是他的姓名从我们的血统传下，}
\verseline{}{皮卢姆努斯是他第四代祖先；而他常以慷慨的}
\verseline{620}{手和众多礼物堆满你的门槛。”}
\verseline{}{高天奥林匹斯之王简短地回答她：}
\verseline{}{“若你只为这将死的青年祈求}
\verseline{}{推迟眼前的死亡，并且明白我只作此让步，}
\verseline{}{就以逃亡带走图尔努斯，从紧逼的命运中夺回他；}
\verseline{625}{我可以宽许到这里。若这些祈求之下}
\verseline{}{藏着更深的恩准，你以为整个战争}
\verseline{}{能够被动摇或改变，那便是以虚妄希望自养。”}
\verseline{}{朱诺流泪说：“你口中不肯给的，若能}
\verseline{}{在心里给我，让这条生命确实留给图尔努斯呢？}
\verseline{630}{如今沉重的结局正等待无辜者，除非}
\verseline{}{我远离了真相。啊，但愿只是虚假的恐惧}
\verseline{}{欺骗了我；愿你既有此力，便把自己的决意转向更好一边！”\footnote{朱庇特讥刺朱诺，朱诺只请求延缓图尔努斯眼前的死亡。朱庇特准她暂时把他从逼近的命运中夺走；朱诺随后仍试探着要求更长久的生命，却没有得到承诺。}}
\verseline{}{说完，她立即从高天投身而下，}
\verseline{}{以雨云束身，驱赶风暴穿过空气，}
\verseline{635}{飞向伊利乌姆人的阵线与劳伦图姆营地。}
\verseline{}{女神从空云中造出一道纤薄、无力的影子，}
\verseline{}{使它具有埃涅阿斯的面容——一种看去令人惊异的幻物——}
\verseline{}{又以达耳达尼亚的兵器装饰，仿造他的盾牌}
\verseline{}{与神圣盔顶的羽冠；她给它空洞的话，}
\verseline{640}{给它没有心智的声音，也摹仿行走的步态，}
\verseline{}{像传说中死后飞荡的形体，}
\verseline{}{或沉睡时戏弄感觉的梦影。}
\verseline{}{幻影在战阵前欣然跳跃，}
\verseline{}{以兵器挑动图尔努斯，又以声音激怒他。}
\verseline{645}{图尔努斯紧逼，从远处掷出呼啸的长枪；}
\verseline{}{幻影背转过去，改变脚步。}
\verseline{}{图尔努斯以为埃涅阿斯转身退走，}
\verseline{}{激荡的心吸入了虚空的希望：}
\verseline{}{“你逃向哪里，埃涅阿斯？不要离开已经许定的婚床；}
\verseline{650}{你渡过海浪寻找的土地，将由我这只右手赐给你。”}
\verseline{}{他高喊着追逐，挥动出鞘的剑，}
\verseline{}{没有看见风正带走他的欢喜。}
\verseline{}{恰有一艘船靠在高耸的岩岸，}
\verseline{}{舷梯已经放下，桥板已经备好；}
\verseline{655}{俄西尼乌斯王曾乘它从克卢西乌姆疆土而来。}
\verseline{}{逃走的埃涅阿斯幻影惊慌地}
\verseline{}{投入其中躲藏；图尔努斯也毫不迟缓，}
\verseline{}{越过阻挡，跃过高高的桥板。}
\verseline{}{他刚触到船首，萨图尔努斯之女便割断缆绳，}
\verseline{660}{使船脱离，卷过回转的海面。}
\verseline{}{轻薄的幻影不再寻找藏身处，}
\verseline{}{而是飞向高处，融入黑云。}
\verseline{}{埃涅阿斯却在战场寻找已经离去的图尔努斯，}
\verseline{}{把许多迎面而来的身体送入死亡；}
\verseline{665}{与此同时，风旋把图尔努斯带到海面中央。\footnote{Conte 2019 将第661—662行移至第664行之后，再接第665行。该版校勘记采用 Brunck 的移行，并说明这样 interea 连接同时发生于两地的事：埃涅阿斯仍在战场作战，图尔努斯同时乘船被带往海上。}}
\verseline{}{他回望，对世事毫不知情，也不感激这次获救；}
\verseline{}{两手一齐伸向星辰，呼喊：}
\verseline{}{“全能的父神，你竟认为我犯下如此重罪，}
\verseline{}{愿我以这样的惩罚偿还吗？}
\verseline{670}{我被带向哪里？从哪里离开？这场逃亡会把我带回何处、带成什么模样？}
\verseline{}{我还能再看见劳伦图姆的城墙与营地吗？}
\verseline{}{那群追随我、追随我兵器的人怎么办？}
\verseline{}{难道我把他们全都——神律所禁啊——留在可耻的死亡中，}
\verseline{}{如今看见他们溃散，听见倒下者的哀号！}
\verseline{675}{我该怎么办？哪一道地缝才足以}
\verseline{}{把我吞没？风啊，不如怜悯我：}
\verseline{}{把船带向峭壁、带向岩石——我图尔努斯心甘情愿地求你们——}
\verseline{}{投入残酷的浅滩与锡耳提斯流沙，}
\verseline{}{让鲁图利人和知晓此事的传闻都不能追上我。”}
\verseline{680}{他说着，心在这一边与那一边摇摆：}
\verseline{}{是为如此大的耻辱发狂，以剑刺进自己，}
\verseline{}{让无情的剑刃穿过胁骨，}
\verseline{}{还是投身海中，游向弯曲的海岸，}
\verseline{}{重新回到特乌克罗斯人的兵刃中。}
\verseline{685}{两条路他都尝试三次，最伟大的朱诺三次}
\verseline{}{拦住他，怜悯青年的心，强行制止。}
\verseline{}{船切过深海，乘顺遂的波浪与潮水滑行，}
\verseline{}{把他送往父亲道努斯的古城。\footnote{本卷第 665—688 行写图尔努斯被朱诺救离战场后的羞耻。他把获救视为抛弃追随者，在自杀与游回战场之间各三次挣扎，又三次被朱诺制止；他甚至祈求船毁于礁石或锡耳提斯流沙，以逃离鲁图利人的目光和“知情的传闻”。}}
\verseline{}{与此同时，墨泽提乌斯遵从朱庇特的意旨，战意燃烧，}
\verseline{690}{接续战斗，袭击正在欢呼的特乌克罗斯人。}
\verseline{}{伊特鲁里亚战阵一齐奔来，所有人都只逼向}
\verseline{}{这一个人，以全部仇恨和密集的兵器逼向这一个人。\footnote{伊特鲁里亚全军把长期仇恨集中在墨泽提乌斯一人身上，背景是第八卷所述的暴政、酷刑、放逐与反抗。海岬和受围野猪的明喻又把他写得坚忍而勇猛。此时的墨泽提乌斯既是遭本国人追杀的流亡者，也是独自迎战全军的强大战士。}}
\verseline{}{他犹如突入浩瀚海面的岩岬，}
\verseline{}{正对狂风，暴露在海中，}
\verseline{695}{承受天空与海洋全部的力量和威胁，}
\verseline{}{自身仍岿然不动；他把多利卡翁之子赫布鲁斯}
\verseline{}{击倒，也击倒拉塔古斯和逃走的帕尔穆斯；}
\verseline{}{他以岩石、以一块巨大的山石碎片}
\verseline{}{迎面砸中拉塔古斯的口与面孔，割断帕尔穆斯的}
\verseline{700}{膝弯，任他无力地翻滚；兵甲则赐给劳苏斯，}
\verseline{}{让他披上肩头，把羽冠插在盔顶。}
\verseline{}{他还杀死佛里吉亚人欧安忒斯，以及帕里斯的同龄}
\verseline{}{伙伴米玛斯；同一夜里，阿米库斯为父，}
\verseline{}{忒阿诺生下米玛斯，怀着火炬的基修斯之女王后}
\verseline{705}{也生下帕里斯；帕里斯长眠故城，\footnote{Conte 2019 校勘记记录本段三项 Bentley 校改：704行 genitore、705行 Paris、710行 pascit。按其正文，忒阿诺为阿米库斯生下米玛斯，基修斯之女王后同夜生下帕里斯；帕里斯死于故城，米玛斯死在劳伦图姆海岸。}}
\verseline{}{劳伦图姆海岸却留住不知情的米玛斯。}
\verseline{}{犹如野猪受猎犬撕咬，从高山上}
\verseline{}{被逐下；多年来，松林遍布的维苏卢斯山}
\verseline{}{庇护它，或劳伦图姆沼泽以芦苇林}
\verseline{710}{养育它；一旦被逼到网中，}
\verseline{}{它站定，凶猛地低吼，鬃毛竖满双肩；}
\verseline{}{没有人有勇气激怒它、走近它，}
\verseline{}{只能从安全的远处以投枪和呼喊逼迫；\footnote{Conte 2019 依 Scaliger 将第 717—718 行移至第 713 行之后，再接第 714—716 行；该版校勘记注明传统传本次序为第 714—718 行。}}
\verseline{}{那头野猪却毫不畏惧，缓缓顾视四方，}
\verseline{}{咬紧牙关，把插在背上的长矛抖落。}
\verseline{}{墨泽提乌斯也正是这样：那些对他怀着正当怒火的人，}
\verseline{715}{没有谁敢拔剑与他近战，}
\verseline{}{只能从远处以投掷兵器和巨大呼声挑衅。}
\verseline{}{阿克戎从古老的科律图斯疆界来到，}
\verseline{720}{他是希腊人，流亡时留下未成的婚礼。}
\verseline{}{墨泽提乌斯远远看见他混战在队伍中央，}
\verseline{}{羽冠紫红，又披着许婚女子的紫衣；}
\verseline{}{犹如饥饿的狮子常在高大畜栏旁游荡，}
\verseline{}{疯狂的饥饿催促它；若碰巧看见}
\verseline{725}{逃走的野山羊，或鹿角高耸的牡鹿，}
\verseline{}{它便欢喜，巨口张开，竖起鬃毛，}
\verseline{}{俯压猎物，紧贴内脏；污黑的血}
\verseline{}{洗过它无度的口——}
\verseline{}{墨泽提乌斯也这样迅猛地冲入密集的敌人。}
\verseline{730}{不幸的阿克戎被击倒，吐出最后的气，}
\verseline{}{脚跟击打黑土，以鲜血染红折断的兵器。}
\verseline{}{奥罗得斯逃走，墨泽提乌斯却不屑于}
\verseline{}{从背后击倒他，也不肯以投枪造成不知来处的伤；}
\verseline{}{他从正面迎上，以勇士对勇士}
\verseline{735}{交锋；不是靠偷袭，而是凭勇武胜过他。}
\verseline{}{他用脚踩住倒下的人，倚着长枪说：}
\verseline{}{“勇士们，高大的奥罗得斯倒下了——战争中不容轻视的一员。”}
\verseline{}{同伴们随着他发出欢喜的凯歌。}
\verseline{}{奥罗得斯却在吐息中说：“胜利者，不论你是谁，}
\verseline{740}{你不能不受报复，也不会欢喜太久；相同的命运}
\verseline{}{也正望着你，不久你将占据同一片土地。”}
\verseline{}{墨泽提乌斯微笑，怒意混在其中：}
\verseline{}{“现在死吧。至于我，留给众神之父、人类之王}
\verseline{}{去裁决。”他说完，从奥罗得斯身体中拔出兵器；}
\verseline{745}{严酷的安息与铁的沉睡压住双眼，}
\verseline{}{眼睛合拢，进入永恒的黑夜。}
\verseline{}{凯狄库斯杀死阿尔卡托乌斯，萨克拉托耳杀死叙达斯佩斯，}
\verseline{}{拉波杀死帕耳忒尼乌斯和力气刚强的俄耳塞斯；}
\verseline{}{墨萨波斯杀死克洛尼乌斯与吕卡翁之子厄里开忒斯：}
\verseline{750}{前者从不受缰绳制驭的马上跌落，躺在地上，}
\verseline{}{后者则徒步作战。吕喀亚人阿吉斯也徒步上前，}
\verseline{}{但继承祖先勇武的瓦勒鲁斯}
\verseline{}{把他击倒；萨利乌斯杀死特罗尼乌斯，涅阿尔刻斯又杀死萨利乌斯——}
\verseline{}{他以标枪和从远处出其不意袭来的箭著称。}
\verseline{755}{玛沃耳斯已让双方的沉重悲哀和交互死亡彼此相等：}
\verseline{}{胜者与败者同样杀人，也同样倒下，}
\verseline{}{这一边不知逃亡，那一边也不知。}
\verseline{}{众神在朱庇特宫中怜悯两边徒然的怒火，}
\verseline{}{怜悯凡人竟有如此巨大的劳苦；}
\verseline{760}{维纳斯从这一边观看，萨图尔努斯之女朱诺从另一边观看，}
\verseline{}{苍白的提西福涅却在万千人中肆虐。\footnote{战斗名录交替写双方杀人者与死者，直到战神使两边的死亡相等。朱庇特宫中的众神怜悯凡人徒劳的怒火，维纳斯与朱诺仍从相反方向观看；第六卷已出现的复仇女神提西福涅在人群中肆虐。}}
\verseline{}{然而墨泽提乌斯挥动巨矛，}
\verseline{}{心潮翻涌地走入平野；犹如巨大的俄里翁，}
\verseline{}{徒步走过涅柔斯水域中央的深处，}
\verseline{765}{劈开道路，双肩耸出波面，}
\verseline{}{或从山巅带回一棵老梣树，}
\verseline{}{踏在地上，头却藏入云层；}
\verseline{}{墨泽提乌斯也这样带着巨大的兵甲逼近。\footnote{俄里翁被写成巨大的涉海者：脚走在涅柔斯的深水中，肩仍高过波面；又如从山巅扛回老梣树，脚在地上、头藏云中。这一明喻把墨泽提乌斯的尺度骤然放大，为他与埃涅阿斯的决斗定标。}}
\verseline{}{埃涅阿斯从长长的队伍中认出他，}
\verseline{770}{准备迎面走去。墨泽提乌斯毫无惊惧，}
\verseline{}{等待心怀壮志的敌人，凭自己的巨躯站定；}
\verseline{}{双眼量好足以投枪的距离：}
\verseline{}{“愿我的右手作我的神，愿我将要掷出的兵器}
\verseline{}{现在临助！劳苏斯，我发誓，要让你披上从那强盗}
\verseline{775}{身体夺来的战利品，使你本人}
\verseline{}{成为战胜埃涅阿斯的纪念。”他说完，从远处投出呼啸的长矛。}
\verseline{}{枪飞去，被盾牌弹开，远远地}
\verseline{}{钉入杰出的安托瑞斯的侧腹，胁骨与髋部之间；}
\verseline{}{安托瑞斯曾是赫尔库勒斯的同伴，从阿尔戈斯被送来，}
\verseline{780}{追随厄万德，定居意大利城中。}
\verseline{}{不幸的人因不属于自己的伤倒下，仰望天空，}
\verseline{}{临死想起可爱的阿尔戈斯。}
\verseline{}{于是虔敬的埃涅阿斯掷出长矛；枪穿过}
\verseline{}{三层青铜拱起的盾面，穿过亚麻衬层，}
\verseline{785}{穿过三张牛皮交织的构造，沉入}
\verseline{}{腹股沟深处，去势也就此耗尽。埃涅阿斯立刻}
\verseline{}{从大腿旁拔剑，因看见伊特鲁里亚人的血而欣喜，}
\verseline{}{炽烈地逼向颤抖的敌人。}
\verseline{}{劳苏斯看见了，因对亲爱父亲的爱}
\verseline{790}{深深叹息，泪水滚过面颊。}
\verseline{}{这里，你的惨死和你最可贵的作为——}
\verseline{}{若后世还肯相信如此伟大的作为——}
\verseline{}{还有你，值得铭记的青年，我都不会略过。}
\verseline{}{墨泽提乌斯向后退步，已无力并受枪伤缠缚；}
\verseline{795}{他一边后退，一边拖着盾上的敌矛；}
\verseline{}{青年冲出，投入两人的兵器之间；}
\verseline{}{埃涅阿斯已经举起右手，正要落下剑击，}
\verseline{}{劳苏斯却来到剑锋之下，以迟延}
\verseline{}{拦住它。同伴们高声呼喊着跟随，}
\verseline{800}{让父亲受儿子的小盾遮护，得以离去；}
\verseline{}{他们投出兵器，以远射}
\verseline{}{击退敌人。埃涅阿斯狂怒，却隐在盾下守住自己。}
\verseline{}{犹如云团有时泻下冰雹，}
\verseline{}{所有耕者都从田野逃开，}
\verseline{805}{所有农人和旅客都藏进安全的避处，}
\verseline{}{或河岸，或高岩的拱穴，}
\verseline{}{直到大雨落尽，他们能在太阳归来后}
\verseline{}{继续白日的工作；埃涅阿斯也这样被四方兵器压住，}
\verseline{}{承受整片战争的雨云，直到雷声退尽，}
\verseline{810}{却还斥责劳苏斯，向劳苏斯发出威胁：}
\verseline{}{“将死之人，你冲向哪里，竟敢做力所不及之事？}
\verseline{}{你的尽责之爱欺骗了毫无戒备的你。”\footnote{埃涅阿斯警告劳苏斯：Fallit te incautum pietas tua。pietas 在这里是促使青年替父承剑的尽责之爱；它“欺骗”毫无戒备的劳苏斯，是因为这种责任把他带向力量不及的死亡。}}
\verseline{}{他仍疯狂地跃动；达耳达尼亚统帅的}
\verseline{}{残酷怒火愈加升起，帕耳卡三女神也为劳苏斯}
\verseline{815}{收取生命最后的丝线：埃涅阿斯把强劲的剑}
\verseline{}{贯入青年的躯干，将整柄剑藏进他体内。}
\verseline{}{剑锋穿过小盾——这威胁者轻薄的兵甲——}
\verseline{}{又穿过母亲以柔金织成的战衣，}
\verseline{}{鲜血涨满衣襟；生命随即悲哀地}
\verseline{820}{穿过空气，走向亡魂，离开躯体。}
\verseline{}{但安喀塞斯之子看见垂死者的面容，}
\verseline{}{看见那张脸以异样的方式转为苍白，}
\verseline{}{便因怜悯深深叹息，伸出右手；}
\verseline{}{父子尽责之爱的形象进入他心中。}
\verseline{825}{“可怜的孩子，因你这些功业，如今}
\verseline{}{虔敬的埃涅阿斯能给你什么，配得上如此禀赋？}
\verseline{}{保留你曾喜爱的兵器吧；若你仍关心这些，}
\verseline{}{我把你送还祖亲的亡魂与骨灰。}
\verseline{}{然而，不幸的孩子，你可以用这一点安慰悲惨的死亡：}
\verseline{830}{你倒在伟大的埃涅阿斯右手之下。”\footnote{埃涅阿斯看见劳苏斯苍白的面容，心中浮现 patriae pietatis imago，即父子尽责之爱的形象。这里的 pietas 承接劳苏斯救父的尽责之爱，也与埃涅阿斯自称 pius 相呼应。他准劳苏斯保留兵甲、归于祖亲的亡魂和骨灰。埃涅阿斯亲自从地上抬起劳苏斯，随后同伴将遗体抬回（831—832、841）。怜悯紧接在他亲手刺杀青年之后。}}
\verseline{}{他又主动斥责迟疑的同伴，亲自从地上抬起}
\verseline{}{劳苏斯；鲜血正污损他依习俗梳理的头发。}
\verseline{}{与此同时，父亲在台伯河水边}
\verseline{}{用清水洗净伤口，使身体得到舒缓，}
\verseline{835}{倚在一棵树的主干上。远处枝头}
\verseline{}{悬着青铜头盔，沉重的兵甲静卧草地。}
\verseline{}{精选的青年站在四周；他负伤喘息，}
\verseline{}{托住颈项，长长的胡须散落胸前；}
\verseline{}{他不断询问劳苏斯，又一再派人}
\verseline{840}{召他回来，传达悲伤父亲的命令。}
\verseline{}{同伴们却流泪把死去的劳苏斯抬在兵甲上，}
\verseline{}{这高大的青年被巨大的创伤击倒。}
\verseline{}{预感灾祸的心远远辨认出哀号；}
\verseline{}{他以厚重尘土污损白发，把双手}
\verseline{845}{伸向天空，又紧紧伏在儿子身上。}
\verseline{}{“儿子，我竟如此贪恋活着，}
\verseline{}{竟容许自己生下的人代我迎向敌人的右手？}
\verseline{}{父亲是靠你的伤口保全的吗，}
\verseline{}{靠你的死亡活下来？唉，直到此刻，毁灭才真正}
\verseline{850}{降临这苦难的我，伤口才真正深入！}
\verseline{}{儿子，也是我以罪行玷污你的名字，}
\verseline{}{因臣民的憎恨被赶下祖传的宝座与权杖。}
\verseline{}{我本应向祖国、向自己人的仇恨付出代价，}
\verseline{}{我本该亲自以千般死法交出这有罪的生命！}
\verseline{855}{如今我仍活着，还没有离开人间与光明，}
\verseline{}{但我会离开。”他说着，以伤腿}
\verseline{}{支起身子；尽管他的力量因深重的伤势而迟滞，}
\verseline{}{他却没有消沉，命人牵来战马。它是他的荣耀，}
\verseline{}{也是他的慰藉；他曾骑着它从}
\verseline{860}{每一场战争胜利归来。他对悲伤的马说：}
\verseline{}{“瑞布斯，若凡人之事能有长久，我们已经活了}
\verseline{}{很久。今日你要么以胜者之姿，带回那些染血的战利品}
\verseline{}{和埃涅阿斯的头颅，同我为劳苏斯的苦痛复仇，}
\verseline{}{要么，若没有力量为我们开路，便与我一同倒下；}
\verseline{865}{最勇猛的战马，我不信你肯忍受}
\verseline{}{外人的命令，屈从特洛伊主人。”\footnote{瑞布斯是墨泽提乌斯有名的战马，曾伴他从每一场战争胜利归来。父亲抱住劳苏斯，承认自己的罪名玷污儿子，也承认臣民对自己的仇恨；随后他把瑞布斯当作共同赴死的伙伴。失子之痛使这位暴君显出深重的悲哀。}}
\verseline{}{他说完，被马背接住，把熟习的肢体安置其上，}
\verseline{}{双手装满锋利的投枪；}
\verseline{}{青铜在头上闪亮，马鬃羽冠森然耸起。}
\verseline{870}{他这样迅疾地冲入人群中央；一颗心中翻涌着}
\verseline{}{巨大的羞耻，狂乱又与哀恸相混。\footnote{Conte 2019 正文由第 871 行直接接第 873 行；该版校勘记注明传统第 872 行在主要传本中缺失。}}
\verseline{}{他三次以巨大呼声召唤埃涅阿斯。}
\verseline{}{埃涅阿斯认出声音，欣然祈求：}
\verseline{875}{“愿众神之父、愿高远的阿波罗使你}
\verseline{}{开始同我正面交手。”}
\verseline{}{他只说这些，便以敌意的长矛迎上。}
\verseline{}{墨泽提乌斯却说：“最残酷的人，你既已夺走我的儿子，}
\verseline{}{还拿什么恐吓我？这正是你能毁灭我的唯一道路。}
\verseline{880}{我不惧死亡，也不向任何神明低头。}
\verseline{}{住口吧；我是来求死的，却先给你}
\verseline{}{带来这些礼物。”他说完，把长矛扭身投向敌人；}
\verseline{}{随后又一支，再一支，绕着巨大的圆周}
\verseline{}{骑马飞驰，金色盾脐却承受一切。}
\verseline{885}{他三次骑马向左绕过站定的埃涅阿斯，}
\verseline{}{一边从手中投枪；特洛伊英雄三次与他一同转动，}
\verseline{}{以青铜盾面带着一整片巨大的矛林。}
\verseline{}{埃涅阿斯厌倦如此拖延、拔出如此多的枪，}
\verseline{}{又被迫以不利的方式接战；}
\verseline{890}{他心中反复思量，终于暴起，}
\verseline{}{把长矛掷进战马两鬓之间的凹陷处。}
\verseline{}{四蹄兽直立起来，以蹄击打空气，}
\verseline{}{随后倒下，压住被抛落的骑者，}
\verseline{}{把他缠在身下，头朝下，以脱出的肩膀重压他。}
\verseline{895}{特洛伊人和拉丁人的呼声燃烧天空。}
\verseline{}{埃涅阿斯飞来，从鞘中拔出剑，}
\verseline{}{站在上方说：“刚烈的墨泽提乌斯如今在哪里，}
\verseline{}{他心中那股狂野的力量又在哪里？”伊特鲁里亚人仰起头，}
\verseline{}{吸入一口天空，恢复神志，回答：}
\verseline{900}{“残酷的敌人，为何斥骂我，以死亡威胁我？}
\verseline{}{杀死我并非神律所禁；我也不是抱着免死之念来到战场，}
\verseline{}{我的劳苏斯更没有为我与你订下这种盟约。}
\verseline{}{只求一件事——若战败的敌人尚可得到宽容：}
\verseline{}{请让泥土覆盖我的身体。我知道，自己的人民怀着}
\verseline{905}{强烈的仇恨围在四周；求你护我免受他们的狂怒，}
\verseline{}{容我进入儿子的坟墓，与他同葬。”\footnote{897—898行仍是埃涅阿斯的问话；随后墨泽提乌斯抬头恢复神志并回答。他已明说自己来求死（881行）。901行 nullum in caede nefas 可解为如此死去不犯神禁，也可解为埃涅阿斯杀死自愿受死的他不犯禁；塞尔维乌斯（10.901）列有这两解。497行剑带图案也用 nefas，所指是婚夜屠杀。临终他只求获葬、免遭本国仇敌毁辱，并与劳苏斯同墓。}}
\verseline{}{他说完，有意以咽喉承受那柄剑，}
\verseline{}{生命随着奔涌的鲜血倾洒在兵甲之上。\footnote{haud inscius 写墨泽提乌斯清楚而主动地以咽喉承受剑；末句又把生命与奔涌的鲜血并写，死亡在动作与流动中完成。}}
